\section{Gluon PDF from Lattice Calculation Using Pseudo-PDF Method}\label{sec:cal-details}

In this work, we use the unpolarized gluon operator defined in Ref.~\cite{Balitsky:2019krf},
%
\begin{equation}\label{eq:gluon_operator}
 {\cal O(z)}\equiv\sum_{i\neq z,t}{\cal O}(F^{ti},F^{ti};z)-\sum_{i,j\neq z,t}{\cal O}(F^{ij},F^{ij};z),
\end{equation}
%
where the operator ${\cal O}(F^{\mu\nu}, F^{\alpha\beta};z) = F^\mu_\nu(z)U(z,0)F^{\alpha}_{\beta}(0)$, $z$ is the Wilson link length, and the field strength $F_{\mu\nu}$ is defined as
%
\begin{equation}\label{}
 F_{\mu\nu} = \frac{i}{8a^2g_0}\left(\mathcal{P}_{[\mu,\nu]}+\mathcal{P}_{[\nu,-\mu]}+\mathcal{P}_{[-\mu,-\nu]}+\mathcal{P}_{[-\nu,\mu]}\right),
\end{equation}
%
where the $a$ is the lattice spacing,
$g_0$ is the strong coupling constant,
and the plaquette $\mathcal{P_{\mu,\nu}}=U_{\mu}(x)U_{\nu}(x+a\hat{\mu})U^{\dag}_{\mu}(x+a\hat{\nu})U^{\dag}_{\nu}(x)$ and $\mathcal{P_{[\mu,\nu]}}=\mathcal{P}_{\mu,\nu}-\mathcal{P}_{\nu,\mu}$.
There is an alternative operator $\sum_{i\neq z,t}{\cal O}(F^{ti},F^{zi};z)$ corresponding to the same matching kernel in Ref.~\cite{Balitsky:2019krf}.
We do not choose this operator, because it vanishes at $P_z=0$ for kinematic reasons, bringing additional difficulties in obtaining the distributions.

Using this operator in Eq.~\ref{eq:gluon_operator}, we calculate lattice gluon matrix elements of the ground-state meson $|0 (P_z) \rangle$ with various boost momenta $P_z$ and Wilson-line displacement lengths $z$.
We then study their dependence on Ioffe time $\nu=zP_z$,
%
\begin{equation}
\mathcal{M}(\nu,z^2) = \langle 0 (P_z)|{\cal O}(z)|0 (P_z)\rangle,
\label{eq:ME_unpol}
\end{equation}
%
calling $\mathcal{M}(\nu,z^2)$ the Ioffe-time distribution (ITD).
To eliminate the ultraviolet divergences in the ITD, we construct the reduced ITD (RITD) by taking the ratio of the ITD to the corresponding $z$-dependent matrix element at $P_z=0$, and further normalize the ratio by the matrix element at $z^2=0$ as done in the first quark pseudo-PDF calculation~\cite{Orginos:2017kos},
%
\begin{equation}
\mathscr{M}(\nu,z^2)=
\frac{\mathcal{M}(zP_z,z^2)/\mathcal{M}(0\cdot P_z,0)}{\mathcal{M}(z\cdot 0,z^2)/\mathcal{M}(0\cdot 0,0)}.
\label{eq:RITD}
\end{equation}
%
The renormalization of ${\cal O}(z)$ and kinematic factors are cancelled out in the RITDs.
The RITD double ratios used here are automatically normalized to one at $z=0$.

The RITDs are related to the pion gluon $g$ and quark $q_S$ PDFs via the pseudo-PDF matching condition~\cite{Balitsky:2019krf}
%
\begin{align}
\mathscr{M}(\nu,z^2)&=\int_0^1 dx \frac{xg(x,\mu^2)}{\langle x \rangle_g}R_{gg}(x\nu,z^2\mu^2) \nonumber \\
&+\frac{P_z}{P_0}\int_0^1 dx \frac{xq_S(x,\mu^2)}{\langle x \rangle_g}R_{gq}(x\nu,z^2\mu^2),
\label{matching-eq}
\end{align}
%
where $\mu$ is the renormalization scale in $\overline{\text{MS}}$ scheme
and $\langle x \rangle_g=\int_0^1 dx \, x g(x,\mu^2)$ is the gluon momentum fraction of the pion.
We can split the gluon-in-gluon $R_{gg}$ contribution in Eq.~\ref{matching-eq} into two parts, which are introduced in Eqs.~21 and 24 in Ref.~\cite{Radyushkin:2018cvn},
%
\begin{align}
&R_{gg}(y,z^2\mu^2)=R_1(y,z^2\mu^2)+R_2(y), \label{matching-kernel}\\
&R_1(y,z^2\mu^2)=-\frac{\alpha_s(\mu)}{2\pi}N_c\ln\left(z^2\mu^2\frac{e^{2\gamma_E+1}}{4}\right)R_B(y), \label{matching-evolve}\\
&R_2(y)=\cos y-\frac{\alpha_s(\mu)}{2\pi}N_c\left( 2R_B(y)+R_L(y)+R_C(y)\right),
\label{matching-scheme}
\end{align}
%
where $R_1(y,z^2\mu^2)$ is the term related to evolution,
$R_2(y)$
is the term related to scheme conversion,
$\alpha_s$ is the strong coupling at scale $\mu$,
$N_c=3$ is the number of colors,
and $\gamma_E=0.5772$ is the Euler-Mascheroni constant.
%
The $z$ in $R_1(y,z^2\mu^2)$ is chosen to be $2e^{-\gamma_E-1/2}/\mu$ so that the log term vanishes, suppressing residuals that contain higher orders of the log term, as discussed in the paper on the one-loop evolution of the pseudo-PDF~\cite{Radyushkin:2018cvn}.
%
The gluon-in-quark kernel $R_{gq}(y,z^2\mu^2)$, along with $R_B(y)$, $R_L(y)$ and $R_C(y)$, are defined in Eqs.~7.21--23 and the paragraph below Eq.~7.23 in Ref.~\cite{Balitsky:2019krf}.

In this work, we first neglect the pion quark PDF, since the total quark PDF is found to be much smaller than the gluon PDF in global fits~\cite{Barry:2018ort,Novikov:2020snp}.
We will later estimate the systematic uncertainty introduced by this assumption.
The gluon evolved ITD (EITD), $G$ is obtained by using the evolution term $R_1(y,z^2\mu^2)$,
%
\begin{align}
G(\nu,\mu,z^2)&=\mathscr{M}(\nu,z^2)\nonumber\\
&+\int^1_0 dx\, R_1(x,z^2\mu^2)\mathscr{M}(x\nu,z^2).
\label{evolution}
\end{align}
%
The $z$ dependence of the EITDs should be compensated by the $\ln{z^2}$ term in the evolution formula.
In principle, the EITD $G$ is free of $z$ dependence and is connected to the lightcone gluon PDF $g(x,\mu^2)$ through the scheme-conversion term $R_2(y)$,
%
\begin{equation}
G(\nu,\mu)=\int_0^1 dx\, \frac{xg(x,\mu^2)}{\langle x \rangle_g} R_2(x\nu),
\label{conversion}
\end{equation}
%
so the gluon PDF $g(x,\mu^2)$ can be extracted by inverting this equation.

On the lattice, we use clover valence fermions on three ensembles with $N_f = 2+1+1$ highly improved staggered quarks (HISQ)~\cite{Follana:2006rc} generated by the MILC Collaboration~\cite{Bazavov:2012xda} with two different lattice spacings ($a\approx 0.12$ and 0.15~fm) and three pion masses (220, 310, 690~MeV).
The masses of the clover quarks are tuned to reproduce the lightest light and strange sea pseudoscalar meson masses used by PNDME Collaboration~\cite{Rajan:2017lxk,Bhattacharya:2015wna,Bhattacharya:2015esa,Bhattacharya:2013ehc}.
%
We use five HYP-smearing~\cite{Hasenfratz:2001hp} steps on the gluon loops to reduce the statistical uncertainties, as studied in Ref.~\cite{Fan:2018dxu}.
We use Gaussian momentum smearing for the quark fields~\cite{Bali:2016lva} to reach higher meson boost momenta.
Table~\ref{table-data} gives the lattice spacing $a$, valence pion mass $M_\pi^\text{val}$ and $\eta_s$ mass $M_{\eta_s}^\text{val}$, lattice size $L^3\times T$, number of configurations $N_\text{cfg}$, number of total two-point correlator measurements $N_\text{meas}^\text{2pt}$, and separation time $t_\text{sep}$ used in the three-point correlator fits for the three ensembles.
This allows us to reach the continuum limit and physical pion mass through extrapolation.
The total amount of measurements vary in $10^5$--$10^6$ for different ensembles.


\begin{table}[!htbp]
\centering
\begin{tabular}{|c|c|c|c|}
\hline
  ensemble  & a12m220 & a12m310 & a15m310 \\
\hline
  $a$ (fm)  & $0.1184(10)$ & $0.1207(11)$  & $0.1510(20)$ \\
\hline
  $M_{\pi}^\text{val}$ (MeV)  & $226.6(3)$ & $311.1(6)$ & $319.1(31)$\\
\hline
  $M_{\eta_s}^\text{val}$ (MeV)  & $696.9(2)$ & $684.1(6)$ & $687.3(13)$\\
\hline
  $L^3\times T$  & $32^3\times 64$ & $24^3\times 64$ & $16^3\times 48$ \\
\hline
  $P_z$ (GeV)  & $[0,2.29]$ &  $[0,2.14]$ &  $[0,2.05]$ \\
\hline
  $N_\text{cfg}$  & 957 & 1013 & 900 \\
\hline
  $N_\text{meas}^\text{2pt}$  &  731,200  & 324,160 & 21,600 \\
\hline
  $t_\text{sep}$  & \{5,6,7,8,9\}   & \{5,6,7,8,9\} & \{4,5,6,7\} \\
\hline
\end{tabular}
\caption{
Lattice spacing $a$, valence pion mass $M_\pi^\text{val}$ and $\eta_s$ mass $M_{\eta_s}^\text{val}$, lattice size $L^3\times T$, number of configurations $N_\text{cfg}$, number of total two-point correlator measurements $N_\text{meas}^\text{2pt}$, and separation times $t_\text{sep}$ used in the three-point correlator fits of $N_f=2+1+1$ clover valence fermions on HISQ ensembles generated by MILC Collaboration and analyzed in this study.
}
\label{table-data}
\end{table}


The two-point correlator for a meson $\Phi$ is
%
\begin{align}\label{}
C_\Phi^\text{2pt}(P_z;t)&=
 \int dy^3 e^{-i y\cdot P_z} \langle \chi_\Phi(\vec{y},t)|\chi_\Phi(\vec{0},0)\rangle \nonumber \\
  &= |A_{\Phi,0}|^2 e^{-E_{\Phi,0}t} + |A_{\Phi,1}|^2 e^{-E_{\Phi,1}t} + ...,
\label{eq:2pt_fit_formula}
\end{align}
%
where $P_z$ is the meson momentum in the $z$-direction, $\chi_\Phi=\bar{q}_1\gamma_5 q_2$ is the pseudoscalar-meson interpolation operator,
$t$ is the Euclidean time,
and $|A_{\Phi,i}|^2$ and $E_{\Phi,i}$ are the amplitude and energy for the ground-state ($i=0$) and the first excited state ($i=1$), respectively.


\begin{figure*}[htbp]
\centering
\centering
\includegraphics[width=0.9\textwidth]{images/Ratio_c957_O0IB_dir-average_p2_z1_Fits.pdf}
\includegraphics[width=0.9\textwidth]{images/Ratio_c957_O0IB_dir-average_p4_z4_Fits.pdf}
\caption{
Example ratio plots (left), one-state fits (second column) and two-sim fits (last 2 columns) from the lightest pion mass $a\approx 0.12$~fm, $M_\pi\approx 220$~MeV for $P_z=2\times 2\pi/L$, $z=1$ (upper row) and $P_z=4\times 2\pi/L$, $z=4$ (lower row).
The gray band shown on all plots is the extracted ground-state matrix element from the two-sim fit using $t_\text{sep}\in[5,9]$.
From left to right, the columns are:
the ratio of the three-point to two-point correlators with the reconstructed fit bands from the two-sim fit using $t_\text{sep}\in [5,9]$, shown as functions of $t-t_\text{sep}/2$,
the one-state fit results for the three-point correlators at each $t_\text{sep}\in[3,9]$,
the two-sim fit results using $t_\text{sep}\in[t_\text{sep}^\text{min},9]$ as functions of $t_\text{sep}^\text{min}$, and
the two-sim fit results using $t_\text{sep}\in[5,t_\text{sep}^\text{max}]$
as functions of $t_\text{sep}^\text{max}$.}
\label{fig:Ratio-fitcomp}
\end{figure*}


The three-point gluon correlators are obtained by combining the gluon loop with pion two-point correlators.
The matrix elements of the gluon operators can be obtained by fitting the three-point correlators to the energy-eigenstate expansion,
%
\begin{align}\label{eq:3ptC}
&C_\Phi^\text{3pt}(z,P_z;t_\text{sep},t) \nonumber \\
&=\int d^3y\, e^{-iy\cdot P_z}\langle \chi_\Phi(\vec y,t_\text{sep})|{\cal O}(z,t)|\chi_\Phi(\vec 0,0)\rangle \nonumber \\
&= |A_{\Phi,0}|^2\langle 0|{\cal O}|0\rangle e^{-E_{\Phi,0}t_\text{sep}} \nonumber \\
&+ |A_{\Phi,0}||A_{\Phi,1}|\langle 0|{\cal O}|1\rangle e^{-E_{\Phi,1}(t_\text{sep}-t)}e^{-E_{\Phi,0}t} \nonumber \\
&+ |A_{\Phi,0}||A_{\Phi,1}|\langle 1|{\cal O}|0\rangle e^{-E_{\Phi,0}(t_\text{sep}-t)}e^{-E_{\Phi,1}t} \nonumber \\
&+ |A_{\Phi,1}|^2\langle 1|{\cal O}|1\rangle e^{-E_{\Phi,1}t_\text{sep}}+ ...,
\end{align}
%
where $t_\text{sep}$ is the source-sink time separation,
and $t$ is the gluon-operator insertion time.
The amplitudes and energies, $A_{\Phi,0}$, $A_{\Phi,1}$, $E_{\Phi,0}$ and $E_{\Phi,1}$, are obtained from the two-state fits of the two-point correlators.
$\langle 0|{\cal O}|0\rangle$, $\langle 0|{\cal O}|1\rangle$ ($\langle 1|{\cal O}|0\rangle$), and $\langle 1|{\cal O}|1\rangle$ are the ground-state matrix element, the ground--excited-state matrix element, and the excited-state matrix element, respectively.
We extract the ground-state matrix element $\langle 0|{\cal O}|0\rangle$ from the two-state fit of the three-point correlators, or a two-state simultaneous ``two-sim'' fit on multiple separation times with the $\langle 0|{\cal O}|0\rangle$, $\langle 0|{\cal O}|1\rangle$ and $\langle 1|{\cal O}|0\rangle$ terms.

To verify that our fitted matrix elements are reliably extracted, we compare to ratios of the three-point to the two-point correlator
%
\begin{equation}\label{eq:Ratio}
R^\text{ratio}(z,P_z;t_\text{sep},t) = \frac{C^\text{3pt}(z,P_z;t_\text{sep},t)}{C^\text{2pt}(P_z;t_\text{sep})};
\end{equation}
%
if there were no excited states, the ratio would be the ground-state matrix element.
The left-hand side of Fig.~\ref{fig:Ratio-fitcomp} shows example ratios for the gluon matrix elements from the lightest pion ensemble, a12m220, at selected momenta $P_z$ and Wilson-line length $z$.
We see the ratios increase with increasing source-sink separation going from 0.60 to 1.08~fm.
At large separation, the ratios begin to converge, indicating the neglect of excited states becomes less problematic.
The gray bands indicate the ground-state matrix elements extracted using the two-sim fit to three-point correlators at five $t_\text{sep}$.
The convergence of the fits that neglect excited states can also be seen in second column of Fig.~\ref{fig:Ratio-fitcomp}, where we compare one-state fits from each source-sink separations:
the one-state fit results increase as $t_\text{sep}$ increases, starting to converge at large $t_\text{sep}$ to the two-sim fit results.

The third and fourth columns of Fig.~\ref{fig:Ratio-fitcomp} show two-sim fits using $t_\text{sep}\in[t_\text{sep}^\text{min},9]$
and
$t_\text{sep}\in[5,t_\text{sep}^\text{max}]$
to study how the two-sim ground-state matrix elements depend on the source-sink separations input into fit.
We observe that the matrix elements are consistent with each other within one standard deviation, showing consistent extraction of the ground-state matrix element, though the statistical errors are larger than those of the one-state fits.
We observe larger fluctuations in the matrix element extractions when small $t_\text{sep}^\text{min}=3$ and 4, or small $t_\text{sep}^\text{max}=6$ and 7, are used. The ground state matrix element extracted from two-sim fits becomes very stable when $t_\text{sep}^\text{min}>4$ and $t_\text{sep}^\text{max}>7$.

Figure~\ref{fig:RITD-fitcomp} shows the RITD of the same examples
$P_z=2\times 2\pi/L$, $z=1$ and $P_z=4\times 2\pi/L$, $z=4$ from two-sim fit results using $t_\text{sep}\in[t_\text{sep}^\text{min},9]$.
The RITD results, which are constructed to suppress lattice fluctuations, are very stable over the range of different fits considered.
For a12m310 and a15m310 ensembles, the $t_\text{sep}$ dependence of RITDs is milder than those from a12m220 ensemble due to the heavier pion mass.
Overall, our ground-state RITDs from the two-sim fit are stable, and we use them to extract the gluon PDF.


\begin{figure}[htbp]
\centering
	\centering
	\includegraphics[width=0.43\textwidth]{images/RITD_c957_O0IB_p2_z1_tsepmin.pdf}
	\includegraphics[width=0.43\textwidth]{images/RITD_c957_O0IB_p4_z4_tsepmin.pdf}
\caption{
Example RITDs from the a12m220 ensemble as functions of $t_\text{sep}^\text{min}$ for $P_z=2\times 2\pi/L$, $z=1$ (top) and $P_z=4\times 2\pi/L$, $z=4$ (bottom).
The two-sim fit RITD results using $t_\text{sep}\in[t_\text{sep}^\text{min},9]$ are consistent with the ones final chosen $t_\text{sep}\in[5,9]$.}
\label{fig:RITD-fitcomp}
\end{figure}
